\documentclass[12pt,a4paper]{article}

\usepackage[utf8]{inputenc}
\usepackage[spanish,mexico]{babel}
\usepackage{booktabs}
\usepackage{tabularx}
\usepackage{geometry}
\geometry{
    a4paper,
    top=3cm,
    bottom=3cm,
    left=3cm,
    right=3cm
}

% Para mejorar la tipografía
\usepackage{mathptmx}
\usepackage{helvet}

\begin{document}

% --- INICIO DE LA PORTADA ---
\begin{titlepage}
    \centering
    \vspace*{2cm}
    
    {\scshape\LARGE Proyecto Integrador \par}
    \vspace{1.5cm}
    
    \rule{\textwidth}{0.4pt}\vspace*{-\baselineskip}\vspace{3.2pt}
    \rule{\textwidth}{1.6pt}\\[\baselineskip]
    {\huge \bfseries Caso Práctico Integrador \par}
    \vspace{0.2cm}
    \rule{\textwidth}{1.6pt}\vspace*{-\baselineskip}\vspace{3.2pt}
    \rule{\textwidth}{0.4pt}
    
    \vspace{2cm}
    
    {\Large\itshape Equipo de Trabajo:\par}
    \vspace{0.8cm}
    {\Large 
        \textbf{Héctor Isaac}\\[0.4cm]
        \textbf{Jaras Sánchez}\\[0.4cm]
        \textbf{Jesel Moreno}
    \par}
    \vfill
    
    {\large \today\par}
\end{titlepage}
% --- FIN DE LA PORTADA ---

\newpage

\section*{Detalles del Caso Práctico Integrador}
\vspace{0.5cm}

% Tabla con el contenido de la imagen
\begin{table}[h!]
    \centering
    \renewcommand{\arraystretch}{1.8} % Da un poco más de espacio entre las filas
    \begin{tabularx}{\textwidth}{>{\bfseries}l X}
        \toprule
        Elemento & Descripción \\
        \midrule
        Tema & Satisfacción laboral en enfermeros de hospitales públicos. \\
        Diseño & No experimental, descriptivo, de corte transversal. \\
        Población & Todos los enfermeros (500) del Hospital General X. \\
        Muestra & 120 enfermeros seleccionados por muestreo aleatorio. \\
        Inclusión & Contrato de tiempo completo, antigüedad mínima de 1 año. \\
        Exclusión & Enfermeros en periodo de vacaciones o incapacidad durante el estudio. \\
        \bottomrule
    \end{tabularx}
\end{table}

\end{document}
